%% ======================================================================
%% Implementation of the  Effect Propagation Process 
%% ======================================================================

\section{The Reference Implementation: DeepCausality}
\label{sec:implementation}



% ============================================================================
% IMPLEMENTATION CHAPTER SKELETON
%
% The chapter mirrors the layered scientific-computing stack introduced in
% the Contribution subsection of the Introduction. Each section describes
% one layer of the stack, in dependency order, with the EPP framework at
% the top and the causal discovery layer above that. The metric crate sits
% horizontally and is treated in its own section within the foundation
% group because it is depended upon by every other layer.
%
% Each section opens with one short sentence stating what the section
% should cover, followed by a TBD marker. Subsections are listed where
% the structure is more granular than a single section can carry.
%
% Prose to be filled in incrementally. The skeleton itself communicates
% that the implementation has structure and that each layer has a
% defensible scope.
% ============================================================================
 
 
\subsection{Chapter Orientation}
\label{sec:implementation_orientation}
 
This section explains how the implementation chapter relates to the framework chapters that precede it, the level of detail the chapter provides, and what the reader gains by reading further.
 
\textbf{[TBD]}
 
 
\subsection{Numerical Foundation}
\label{sec:implementation_num}
 
This section describes the \texttt{deep\_causality\_num} crate, which provides the algebra trait hierarchy, Float106 high-precision arithmetic, and the Complex, Quaternion, and Octonion division algebras that the rest of the stack builds on.
 
\subsubsection{Algebra Trait Hierarchy}
The trait hierarchy from Magma through DivisionAlgebra, with the marker traits (Associative, Commutative, Distributive) that differentiate intermediate structures, and the implementation strategy that encodes mathematical laws as compile-time promises.
 
\textbf{[TBD]}
 
\subsubsection{Float106 Double-Double Precision}
The double-double precision arithmetic implementation, the comparison with IEEE binary128, and the design choices that allow Float106 to be available on stable Rust while delivering competitive precision and speed.
 
\textbf{[TBD]}
 
\subsubsection{Complex, Quaternion, and Octonion}
The division algebra implementations, their classification within the algebra trait hierarchy (Field, AssociativeRing, DivisionAlgebra respectively), and the Rotation trait that supports geometric operations.
 
\textbf{[TBD]}
 
\subsubsection{No-Std Support and Design Constraints}
The macro-free, unsafe-free, and dependency-free implementation strategy, the no-std support via libm, and the rationale for these constraints.
 
\textbf{[TBD]}
 
 
\subsection{Higher-Kinded Types}
\label{sec:implementation_haft}
 
This section describes the \texttt{deep\_causality\_haft} crate, which provides higher-kinded types in Rust through type-level encoding and the witness pattern, and supplies the Effect, Functor, Applicative, Monad, and CoMonad trait families that all higher layers consume.
 
\subsubsection{The Witness Pattern}
The technique for encoding higher-kinded types in Rust via witness types, the role of the HKT trait family, and the trade-offs compared to languages with native HKT support.
 
\textbf{[TBD]}
 
\subsubsection{Algebraic Traits}
The Effect, Functor, Applicative, Monad, CoMonad, and BoundedComonad trait families, with their laws and their relationships to one another.
 
\textbf{[TBD]}
 
\subsubsection{Arity-Five Effect Types}
The Effect5 trait family with its five type parameters (value, state, context, error, log) and the rationale for arity five as the right level of expressiveness for the EPP's monadic effect system.
 
\textbf{[TBD]}
 
 
\subsection{Metric Specification}
\label{sec:implementation_metric}
 
This section describes the \texttt{deep\_causality\_metric} crate, which provides the single source of truth for metric signatures across the stack and eliminates the class of consistency bugs that would otherwise arise when general-relativity and particle-physics conventions are implemented independently.
 
\subsubsection{Sign Convention Types}
The East Coast and West Coast convention types, the Lorentzian metric trait, and the type aliases (RelativityMetric, ParticleMetric) that document intent.
 
\textbf{[TBD]}
 
\subsubsection{Clifford Signature Space}
The $Cl(p, q, r)$ signature representation, the Metric enum and its variants (Euclidean, NonEuclidean, Minkowski, PGA, Generic, Custom), and the convention conversion routines.
 
\textbf{[TBD]}
 
\subsubsection{Cross-Crate Consumption}
How the multivector, topology, and physics crates consume the metric specification, and how horizontal placement of the metric crate enforces consistency without manual coordination.
 
\textbf{[TBD]}
 
 
\subsection{Tensor Layer}
\label{sec:implementation_tensor}
 
This section describes the \texttt{deep\_causality\_tensor} crate, which provides multi-dimensional arrays with stride-based memory layout, broadcasting, Einstein summation, and higher-kinded type instances.
 
\subsubsection{Stride-Based Memory Layout}
The flat-vector backing store with stride-based indexing, the memory-layout decisions that support broadcasting without copying, and the performance characteristics of element access and shape manipulation.
 
\textbf{[TBD]}
 
\subsubsection{Broadcasting and Element-Wise Operations}
The broadcasting rules and their alignment with established conventions, the binary operation engine, and the support for tensor-scalar and tensor-tensor arithmetic.
 
\textbf{[TBD]}
 
\subsubsection{Einstein Summation}
The EinSumOp abstraction, the supported operations (matrix multiplication, dot product, trace, batch matrix multiplication, element-wise product), and the AST-based execution strategy.
 
\textbf{[TBD]}
 
\subsubsection{Higher-Kinded Type Instances}
The Functor, Applicative, Monad, and CoMonad instances for CausalTensor, with examples of how the instances enable composition with other layers and with effect systems.
 
\textbf{[TBD]}
 
 
\subsection{Multivector Layer}
\label{sec:implementation_multivector}
 
This section describes the \texttt{deep\_causality\_multivector} crate, which provides Clifford algebras over the dynamic signature space with pre-configured constructors for the algebras commonly used in physics and geometry.
 
\subsubsection{Universal Multivector}
The CausalMultiVector type that represents scalars, vectors, bivectors, and higher-grade blades within a single type, the geometric product implementation, and the grade projection operations.
 
\textbf{[TBD]}
 
\subsubsection{Pre-Configured Algebras}
The pre-configured constructors for Pauli algebra, spacetime algebra in both conventions, conformal geometric algebra, projective geometric algebra in three dimensions, the Dixon algebra, and the Grand Unified Algebra hosting $\mathfrak{spin}(10)$ symmetry.
 
\textbf{[TBD]}
 
\subsubsection{Higher-Kinded Type Instances}
The Functor, Applicative, and Monad instances, with the bind operation realizing the tensor product of algebras and the consequences for dimension-changing compositions.
 
\textbf{[TBD]}
 
\subsubsection{Quantum State Vectors}
The HilbertState type for quantum state representations, the gate operations, and the integration with the multivector machinery.
 
\textbf{[TBD]}
 
 
\subsection{Topology Layer}
\label{sec:implementation_topology}
 
This section describes the \texttt{deep\_causality\_topology} crate, which provides graphs, hypergraphs, simplicial complexes, manifolds, and point clouds together with differential geometry operators and a lattice gauge field framework.
 
\subsubsection{Topological Types}
The Graph, Hypergraph, SimplicialComplex, Manifold, PointCloud, Lattice, CellComplex, Chain, Skeleton, Simplex, and Topology types, with their mathematical correspondences and their typical uses.
 
\textbf{[TBD]}
 
\subsubsection{Topological Algorithms}
The Vietoris-Rips triangulation, the Euler characteristic computation, the boundary and coboundary operators, and the chain algebra operations.
 
\textbf{[TBD]}
 
\subsubsection{Differential Geometry}
The exterior derivative, Hodge star, codifferential, and Hodge-Laplacian operators on discrete differential forms, with the Hodge theory machinery for detecting topological features.
 
\textbf{[TBD]}
 
\subsubsection{Lattice Gauge Field Theory}
The GaugeField, LatticeGaugeField, and LinkVariable types, the supported gauge groups (U(1), SU(2), SU(3), Lorentz), and the verification against twenty-four reference results from Creutz's lattice gauge theory text.
 
\textbf{[TBD]}
 
\subsubsection{Higher-Kinded Type Instances}
The Functor, BoundedComonad, and BoundedAdjunction instances, with the comonadic extend operation as the foundation for graph convolutions and cellular automata.
 
\textbf{[TBD]}
 
 
\subsection{Physics Layer}
\label{sec:implementation_physics}
 
This section describes the \texttt{deep\_causality\_physics} crate, which builds on the four mathematical layers below it to provide physics kernels organized into thematic domains, propagating-effect wrappers, and full physical theories implemented as gauge theories.
 
\subsubsection{Kernel-Versus-Theory Separation}
The architectural separation between pure stateless kernels for isolated equations and full physical theories that use gauge fields and geometric algebra, and the rationale for keeping both levels of abstraction available.
 
\textbf{[TBD]}
 
\subsubsection{Physics Domains}
The thematic domain organization (astrophysics, condensed matter, classical dynamics, electromagnetism, fluid dynamics, materials, magnetohydrodynamics, nuclear physics, photonics, quantum mechanics, relativity, thermodynamics, waves) and the rationale for the chosen scope.
 
\textbf{[TBD]}
 
\subsubsection{Propagating-Effect Wrappers}
The wrapper functions that lift each kernel into the PropagatingEffect monad, the design pattern that allows kernels to be embedded in causal effect chains, and the implications for composition with the EPP framework.
 
\textbf{[TBD]}
 
\subsubsection{Gauge Theories}
The implementation of general relativity, electromagnetism, weak force, and electroweak unification as gauge theories over the topology crate's gauge field abstraction, with the gauge group as the only point of difference between theories.
 
\textbf{[TBD]}
 
\subsubsection{Type-Safe Units and Generic Precision}
The newtype wrappers for physical quantities, the prevention of dimensional errors at the type level, and the genericity over floating-point type that allows simulations to be run at single, double, or Float106 precision.
 
\textbf{[TBD]}
 
 
\subsection{Causal Layer}
\label{sec:implementation_causal}
 
This section describes the four crates that constitute the EPP framework: the core types, the main framework, the UltraGraph hypergraph data structure, and the Uncertain type for first-order probabilistic data.
 
\subsubsection{Core Types and Effect System}
The \texttt{deep\_causality\_core} crate's PropagatingEffect, PropagatingProcess, CausalMonad with the bind operation, the EffectValue type, the EffectLog audit trail, and the Causal Effect System runtime.
 
\textbf{[TBD]}
 
\subsubsection{The ControlFlowBuilder}
The compile-time-verified static execution graph capability, the type-safe topology enforcement, the zero-allocation execution loop, and the strict-zst feature for safety-critical and hard-real-time deployment.
 
\textbf{[TBD]}
 
\subsubsection{The Causaloid}
The Causaloid type, its three isomorphic forms (Singleton, Collection, Graph), the AggregateLogic for collections, the context-aware variants, and the implementation of the MonadicCausable trait.
 
\textbf{[TBD]}
 
\subsubsection{Context}
The Context structure, the contextoid nodes, the Euclidean and non-Euclidean representations, the linear and non-linear temporal structures, and the static and dynamic context modes.
 
\textbf{[TBD]}
 
\subsubsection{Causal State Machine}
The CSM that links causal inference to deterministic action, the state representation, the action dispatch mechanism, and the integration with the Effect Ethos.
 
\textbf{[TBD]}
 
\subsubsection{Teloid and Effect Ethos}
The Teloid as a single unit of intent, the Effect Ethos that aggregates teloids into a rule set, the defeasible deontic calculus that resolves normative conflicts, and the integration with the Causal State Machine.
 
\textbf{[TBD]}
 
\subsubsection{UltraGraph}
The \texttt{ultragraph} crate's two-phase hypergraph data structure, the design choices that support performance-constrained hypergraph composition, and the use within the CausaloidGraph and the Context Hypergraph.
 
\textbf{[TBD]}
 
\subsubsection{Uncertain Type}
The \texttt{deep\_causality\_uncertain} crate's implementation of Bornholt's Uncertain<T>, the support for full probability distributions and computation graphs, and the lifting of deterministic and probabilistic effects into the uncertain distribution.
 
\textbf{[TBD]}
 
 
\subsection{Causal Discovery Layer}
\label{sec:implementation_discovery}
 
This section describes the \texttt{deep\_causality\_discovery} crate, which sits at the top of the stack and bridges from raw observational data to executable causal models through a typestate-encoded pipeline.
 
\subsubsection{Typestate Pipeline}
The Rust typestate pattern as applied to the discovery pipeline, the sequence of stages (configuration, data loading, cleaning, feature selection, causal discovery, analysis, finalization), and the compile-time enforcement of valid stage transitions.
 
\textbf{[TBD]}
 
\subsubsection{Data Loading and Cleaning}
The CSV and Parquet data loaders, the data-cleaning interfaces, and the design choices for handling messy observational data.
 
\textbf{[TBD]}
 
\subsubsection{Feature Selection}
The minimum-redundancy maximum-relevance feature selection implementation and its role in preparing data for the discovery primitive.
 
\textbf{[TBD]}
 
\subsubsection{The SURD-States Algorithm}
The decomposition of causal influences into Synergistic, Unique, and Redundant components, the algorithmic structure, and the interpretation of each component class.
 
\textbf{[TBD]}
 
\subsubsection{From Discovery to Causal Model}
The translation from SURD analysis to CausaloidGraph construction, the mapping from synergistic and unique influences to graph topology, and the use of redundant influences in AggregateLogic decisions.
 
\textbf{[TBD]}
 
 
\subsection{Cross-Layer Composition}
\label{sec:implementation_composition}
 
This closing section walks through how the layers compose in practice, with the avionics example demonstrating the dual-architecture composition within the causal layer and the chronogravimetric example demonstrating composition across the entire stack.
 
\textbf{[TBD]}
 
%% ======================================================================
%% Validation and verification 
%% ======================================================================

\subsection{Validation and Verification}
\label{sec:implementation_validation}

\subsubsection{Overview}

Regulated industries such as avionics, robotics, or defense have to comply with strict regulatory requirements for software certification. Rust as a language provides a significant amount of safety relevant features that are difficult to achieve with C/C++ while still providing compatibility with existing code C/C++. Furthermore, the existence of a certified Rust compiler enables the validation and verification process of Rust systems in regulated industries. Build on this foundation, the DeepCausality project, while not certified, provides a plausible pathway towards certification. 

\subsubsection{Toolchain}

Rust a a language was chose because of its memory safety and the availability of certified compilers from vendors such
as Ferrocene\footnote{\url{https://ferrocene.dev/en}} (ISO 26262 ASIL D, IEC 61508 SIL 3, IEC 62304 / Class C) HighTec EDV-Systeme\footnote{\url{https://hightec-rt.com/products/rust-development-platform}} (ISO 26262 ASIL D), and AdaCore\footnote{\url{https://www.adacore.com/gnatpro-rust}}. 

\subsubsection{Hermetic Builds and Tests}

Beyond the language-level assurances, DeepCausality implements a robust build and test infrastructure specifically designed for high-integrity environments. Regulatory standards like ISO 26262 (Automotive Safety), IEC 61508 (Functional Safety), DO-178C (Avionics Software), and IEC 62304 (Medical Device Software) demand that the software build process fulfils:

\begin{itemize}
	\item Reproducibility: Given the same inputs, the system must always produce the same outputs. 
	\item Traceability: Every output (including test results) must be traceable back to its exact inputs.
	\item Independence: Tests must be independent of each other. The failure of one component should not mask or cause the failure of another.
	\item Controlled Environment: The testing environment must be fully specified and controlled.
\end{itemize}


The entire DeepCausality mono-repository is built and tested using the Bazel build system. Bazel's core principle of running tests in fully isolated, hermetic sandboxes ensures that each test executes independently, free from external state or interference from other tests. Bazel enforces hermeticity by:

\begin{itemize}
	\item Reproducible Inputs: Bazel requires all inputs (source code, compilers, libraries, tools) to be explicitly declared. This makes builds and tests fully reproducible.
	\item Isolation: Each test runs in its own isolated environment (a "sandbox"). This guarantees isolation and prevents side effects.
	\item Traceability: Each artifact in Bazel can be fully traced via the build graph. Bazel supports search and through the build graph which means one can quickly find all affected dependencies in case an issue has been found. 
\end{itemize}

\subsubsection{Design Principles}

The DeepCausality project has been designed with regulatory requirements in mind to streamline the validation and verification process that precedes certification. Specifically, DeepCausality is built upon three core principles that establish a baseline for trustworthiness:

\begin{itemize}
	\item Zero External Dependencies
	\item Zero Unsafe Code
	\item Zero Macros in Libraries
\end{itemize}

\subsubsection{Zero External Dependencies}

The DeepCausality codebase, including its custom tensor and random number generation primitives, compiles without external third-party dependencies by default. The decision to implement a custom CausalTensor as foundation for the various algorithms used in the DeepCausality project as well as the decision to implement a custom random number generation that underpins the Uncertain crate fundamentally enabled the zero dependency achievement. Among other benefits, this ensures that the Effect Ethos's deontic inference process operates with predictable and auditable determinism, which is critical for proofing end to end alignment. 
The zero external dependency principle reduces supply chain risk by eliminating unknown code from external sources and effective mitigates a key security concern in high-integrity systems. It ensures that the entire intellectual supply chain is auditable and simplifies cross-platform compilation for diverse embedded hardware targets.


The only exception is an optional feature flag, os-random, in the deep\_causality\_rand crate. When enabled, this flag introduces a dependency on a thin Rust wrapper for libc to access the operating system’s secure random number generator. This is clearly documented and is strictly opt-in for developers who require a cryptographically secure random numbers provided by the host system. Alternatively, the deep\_causality\_rand crate also contains a flexible trait that exposes the random number generator (RNG) functionality. For regulated industries, the trait-based architecture for RNG allows them to seamlessly swap in their own certified hardware RNG binding by simply implementing one trait, which means DeepCausality becomes compatible with  existing certified hardware. 

For the development process, there are a handful of external dependencies i.e. Criterion for running benchmarks or CSV for reading data files, but these are very limited, all contained outside the DeepCausality codebase, and only applicable to the development process with the understanding that none of these will be part of a production build.  

\subsubsection{Zero Unsafe Code}

The entire DeepCausality codebase including all internal dependencies, examples, and tests, are implemented entirely in safe Rust. This preserves Rust's guarantee of memory safety and therefore eliminating an entire class of catastrophic failures (e.g., buffer overflows, data races) at compile time. It significantly reduces the burden of memory safety audits, which are typically one of the most complex and costly aspects of high-integrity software.

\subsubsection{Zero Macros in Libraries}

The entire DeepCausality codebase is free of macros to ensures maximum code clarity and audibility. The entire codebase can be read, understood, and analyzed line-by-line, which is essential for formal verification processes and human review in regulated environments. The only area where macros had to be used is in bulk testing for generic traits as, for example, in the deep\_causality\_num trait. The reason is quite obvious, core numerical traits such as "mul" that defines the multiplication operation for every numerical type would require identical tests for every possible numerical type in Rust. In those cases, a macro is used to derive the test mainly to keep the test files largely free of duplicates. It is important to understand that this practice is contained only to tests and only applies to very specific foundational crates i.e. deep\_causality\_rand and deep\_causality\_num. 


These key principles are complemented by a rigorous testing regimen that maintains a 3-month rolling average code coverage of 97\% across its 80,000+ lines of Rust code. While code coverage does not constitute formal verification, it is a critical quantitative indicator of testing thoroughness. For certification bodies, this high coverage, coupled with the inherent safety of Rust and the absence of external dependencies, substantially lowers the time, effort and cost associated with verification and validation activities. Combined, these three key principles enable a number of properties particular relevant to regulated industries:

\begin{itemize}
	\item Supply Chain Security: By being entirely self-contained, DeepCausality guarantees the integrity of its code. 
	\item Cross-Compilation \& Embedded Systems: The std-only approach makes DeepCausality inherently portable across a vast array of embedded target platforms supported by Rust. 
	\item Performance \& Development Velocity: Because the DeepCausality codebase is self-contained, the entire mono-repository (est. 80K Lines of Code), completes a full clean build within 1o seconds translates directly to developer productivity. 
\end{itemize}

The DeepCausality project provides a clear and plausible pathway towards regulatory certification because of the existence of a certified Rust compiler, the already established hermetic build system, and the achieved zero dependencies, unsafe, and macros best practice across the entire project codebase. Systems build with DeepCausality that seek certification still need to comply with all regulatory requirements, but because of the careful preparations already made, the process of validation and verification becomes manageable for an experienced team.


\newpage
